% D.32: reviewed route obstruction; no universal section claim.
\subsection{Full modular levels need not increase rational unit rank}

\begin{claim}{D.32}{Proved route obstruction}\label{clm:D32}
Let $s$ be a continuous section of the rational tripod
$U=\mathbf P^1_{\mathbf Q}\setminus\{0,1,\infty\}$, let
$\mathcal E_\lambda:y^2=x(x-1)(x-\lambda)$ be the Legendre family,
and put $W_N=s^*\mathcal E[N]$ for an odd integer $N\ge3$.
The symplectic-isomorphism cover
\[
 Y_{s,N}=\operatorname{Isom}_{\mathrm{sympl}}
               (W_N,\mathcal E[N])\longrightarrow U
\]
is a geometrically connected finite \'{e}tale neighbourhood of $s$.
If
\[
 \rho_{s,N}(G_{\mathbf Q})=
                 \operatorname{GL}_2(\mathbf Z/N),
 \tag{D.32.1}
\]
its smooth completion has exactly three closed boundary points over
$\mathbf Q$, one above each original cusp label. In particular
\[
 \operatorname{rank}
       \bigl(\mathcal O(Y_{s,N})^\times/\mathbf Q^\times\bigr)=2.
 \tag{D.32.2}
\]
Every rational unit $v$ satisfies
\[
          v^m=q\,\lambda^a(1-\lambda)^b
 \quad\text{for some }m>0,\ a,b\in\mathbf Z,\
 q\in\mathbf Q^\times.
 \tag{D.32.3}
\]
Thus a unit regular and nonzero at cusps above two different
original labels is constant. The same conclusions hold after
quotienting the cover by its central sign.

There is an explicit infinite tower satisfying
\textup{(D.32.1)}: take the point section at $\lambda=3$ and
$N=5^r$, $r\ge1$. Its degrees, genera and geometric cusp counts
grow without bound, while the rational unit rank remains two.

Consequently these levels do not automatically furnish the
unit-rich neighbourhood required in D.31, and in the
maximal-image case they cannot supply its unit when two different
cusp labels occur. This is a limitation of this route. It is not
a nongeometric-section example or a conclusion about every finite
\'{e}tale neighbourhood.
\end{claim}

\begin{proof}
\proofparagraph{The actual integral cover and its geometric group.}
The Weil pairing makes $W_N$ and $\mathcal E[N]$ symplectic
finite \'{e}tale modules with pairing in $\mu_N$.
After an \'{e}tale trivialization their symplectic-isomorphism
space is the constant finite set
$G_N=\operatorname{SL}_2(\mathbf Z/N)$.
This proves representability and \'{e}taleness; it is also the
level-isomorphism construction in
\citep[\S17, Proposition~17.6]{D_ConradModularNotes}.
Pullback by $s$ gives $\operatorname{Isom}_{\mathrm{sympl}}(W_N,W_N)$.
Its identity is Galois fixed, so $s$ lifts.
In particular no comparison of rational Tate modules is being
replaced by an unproved integral-lattice comparison.

The two geometric Legendre monodromies are, in a symplectic basis,
\[
 A=\begin{pmatrix}1&2\\0&1\end{pmatrix},\qquad
 B=\begin{pmatrix}1&0\\-2&1\end{pmatrix}.
 \tag{D.32.4}
\]
They follow directly from the two Picard--Lefschetz squares
\citep[\S1.1, pp.~20--21]{D_CarlsonMullerStachPeters2017}.
Since $2$ is invertible modulo $N$, their powers contain every
upper and lower elementary unipotent.
These generate $G_N$: over a local factor of $\mathbf Z/N$,
an entry of a unimodular column is a unit, so elementary row
operations reduce it to the first standard vector, after which
the remaining upper entry is removed. The Chinese remainder
theorem combines the elementary operations over all local factors.
Thus the cover is geometrically connected and has geometric
deck group $G_N$.

\proofparagraph{The full and projective cusp fibres.}
The third geometric monodromy is
\[
 C=(AB)^{-1}=\begin{pmatrix}1&-2\\2&-3\end{pmatrix}.
 \tag{D.32.5}
\]
The first two inertia groups have order $N$ and are conjugate
to the full upper unipotent group $U_N$.
Since $C=-I+M$, $M^2=0$, and $M$ has a unit entry modulo $N$,
its order is $2N$. Its inertia group is conjugate to
$\{\pm I\}U_N$.
With $I_c$ denoting these inertia subgroups, the geometric
cusp fibres are the left coset sets $I_c\backslash G_N$.
Their cardinalities are therefore
\[
 \frac{|G_N|}{N},\quad
 \frac{|G_N|}{N},\quad
 \frac{|G_N|}{2N},
 \qquad
 |G_N|=N^3\prod_{\ell\mid N}(1-\ell^{-2}).
 \tag{D.32.6}
\]
The central-sign quotient has $|G_N|/(2N)$ cusps over
each label. It is geometrically the congruence cover
$X(2N)\to X(2)$, by combining the specified level-two
structure and the level-$N$ structure modulo the elliptic
automorphism $-1$. The full torsion cover has twice its degree;
these two covers have not been identified.
For the usual primitive-vector and modular-function-field
descriptions, see
\citep[Lemma~2.1 and Proposition~2.2, pp.~1206--1207]
{D_SutherlandZywina2017}.

\proofparagraph{Arithmetic transitivity on each cusp fibre.}
Choose a tangential splitting at $c$, a local Tate basis
and a basis of $W_N$. A symplectic level isomorphism is
represented by $h\in G_N$, with Galois action
\[
 h\longmapsto
       \rho_c(\sigma)h\rho_{s,N}(\sigma)^{-1}.
 \tag{D.32.7}
\]
The local Legendre degeneration is potentially multiplicative.
After its quadratic splitting twist, Tate uniformization gives
the filtration with characters $\chi_N$ and $1$; its
Galois equivariance is recalled in
\citep[\S19, p.~44]{D_ConradModularNotes}.
Restrict the quadratic splitting character to the chosen
tangential copy of $G_{\mathbf Q}$, calling it $\delta_c$.
For $\sigma\in G_{\mathbf Q(\mu_N)}$ this gives
\[
            \rho_c(\sigma)=\delta_c(\sigma)u_\sigma,
                   \qquad u_\sigma\in I_c.
 \tag{D.32.8}
\]
At infinity the horizontal sign already lies in $I_\infty$.
Thus the same formula applies there; nonsplit tangential
degenerations have not been excluded.

Passing \textup{(D.32.7)} to left cosets, and using that
$\delta_c$ is scalar, gives the right action
\[
 I_ch\longmapsto
       I_ch\bigl(\delta_c(\sigma)\rho_{s,N}(\sigma)\bigr)^{-1}.
 \tag{D.32.9}
\]
The determinant of $\rho_{s,N}$ is $\chi_N$, so
\textup{(D.32.1)} implies that its restriction to
$G_{\mathbf Q(\mu_N)}$ maps onto $G_N$.
Let $H_c$ be the image of its quadratic twist in
\textup{(D.32.9)}. It has full image in
$G_N/\{\pm I\}$; hence $H_c\{\pm I\}=G_N$ and
$[G_N:H_c]\le2$.
But $G_N$ has no quotient of order two: its elementary
unipotent generators all have odd order dividing $N$.
A proper subgroup of index two would give such a quotient.
Therefore $H_c=G_N$, which acts transitively on
$I_c\backslash G_N$.

The three original labels are distinct rational points.
They cannot be interchanged by Galois, and each of their
fibres is now a single orbit. This proves that there are
exactly three closed boundary points over $\mathbf Q$.
Passing to the sign quotient preserves the same conclusion.

\proofparagraph{The rational-unit calculation.}
Let $D_0,D_1,D_\infty$ be the closed boundary divisors,
with coefficient one on each geometric point in the
corresponding orbit. Geometric connectedness gives an injection
\[
 \mathcal O(Y_{s,N})^\times/\mathbf Q^\times
 \lhook\joinrel\longrightarrow
 \operatorname{Div}^0_{\partial C_{s,N}}
                           (C_{s,N})(\mathbf Q).
 \tag{D.32.10}
\]
The latter is the kernel of the degree map from the free
rank-three group on the $D_c$, so has rank two.
The pullbacks of $\lambda$ and $1-\lambda$ are independent,
proving \textup{(D.32.2)}.

More explicitly, writing $e_c$ for the ramification indices,
\[
 \begin{split}
 \operatorname{div}(\lambda)&=e_0D_0-e_\infty D_\infty,\\
 \operatorname{div}(1-\lambda)&=e_1D_1-e_\infty D_\infty.
 \end{split}
 \tag{D.32.11}
\]
These span the degree-zero boundary space over $\mathbf Q$.
For any $v$, clear the denominators expressing
$\operatorname{div}(v)$ in this basis. The quotient of the
resulting functions has zero divisor and is a rational constant,
proving \textup{(D.32.3)}.
If $v$ has order zero at cusps over labels $0,1$, then
$a=b=0$. For labels $0,\infty$ one gets $a=0,a+b=0$;
for labels $1,\infty$ one gets $b=0,a+b=0$.
In each case $v$ is constant.
All arguments apply verbatim to the sign quotient.

\proofparagraph{The explicit maximal-image tower.}
For $E_3:y^2=x(x-1)(x-3)$ one computes
\[
                     \Delta=576,\qquad c_4=112.
 \tag{D.32.12}
\]
At the local prime $3$, the model is minimal with
multiplicative reduction of Kodaira type $I_2$:
$v_3(c_4)=0$ and $v_3(\Delta)=2$.
After at most an unramified quadratic extension its Tate
parameter has valuation two. On $T_5(E_3)$, inertia therefore
contains the whole subgroup
\[
 U(\mathbf Z_5)=
    \left\{\begin{pmatrix}1&z\\0&1\end{pmatrix}:z\in\mathbf Z_5\right\}.
 \tag{D.32.13}
\]
Indeed the inertial Kummer parameter is $2t_5$, which is
surjective because $2$ is a $5$-adic unit; a nonsplit
multiplicative twist is unramified and does not change it.

Over $\mathbf F_7$ the only affine points are
$(0,0),(1,0),(3,0)$. Thus $\#E_3(\mathbf F_7)=4$,
and a Frobenius $g$ has characteristic polynomial
$T^2-4T+7$. Its discriminant modulo five is the nonsquare
$3$. If $e_1$ spans the fixed line of
\textup{(D.32.13)}, the vectors $e_1,e_2=ge_1$ form a
$\mathbf Z_5$-basis. In this basis,
\[
          g=\begin{pmatrix}0&-7\\1&4\end{pmatrix},
            \qquad gE_{12}g^{-1}=-\frac17E_{21}.
 \tag{D.32.14}
\]
The original upper unipotent group remains full in this
basis, and its conjugate is the full lower unipotent group.
They generate $\operatorname{SL}_2(\mathbf Z_5)$.
The determinant is the surjective $5$-adic cyclotomic
character, so the image is
$\operatorname{GL}_2(\mathbf Z_5)$.
Every level $5^r$ therefore satisfies
\textup{(D.32.1)}.

Formula \textup{(D.32.6)} proves unbounded degree and cusp
count, and Riemann--Hurwitz proves unbounded genus.
At level five the full cover has degree $120$, cusp counts
$24,24,12$ and genus $31$; its sign quotient has degree
$60$, cusp counts $12,12,12$ and genus $13$.
The section used for this tower is an actual point section.
It demonstrates failure of automatic unit-rank growth,
not failure of the section conjecture.
\end{proof}

\proofparagraph{Why Manin--Drinfeld does not repair the rank gap.}
For a $\mathbf Q$-form of a congruence modular curve with
its cuspidal boundary, let $c$ be the number of closed cusps.
The rational unit rank is $c-1$.
To justify the lower bound, any rational degree-zero
cuspidal divisor is torsion in the geometric Picard group
by Drinfeld's original Theorem~1
\citep[p.~83, Russian original]{D_DrinfeldCusps1973}.
Some multiple is principal over a finite Galois extension;
the norm of that function makes a further positive multiple
principal over $\mathbf Q$.
Thus principal rational boundary divisors have finite index
in the rank-$(c-1)$ degree-zero lattice.
The corresponding cusp description is also recalled in
\citep[\S\S2.3--2.4, Lemma~2.3]{D_ZywinaRunge}.
This gives rank two for our projectivized modular curves.
The full-cover proof above already obtained rank two from
the two base units and needed no Manin--Drinfeld input.
A growing number of geometric cusps inside three Galois
orbits supplies no extra rational directions.

\proofparagraph{A conditional bound for an open torsion image.}
Fix an odd prime $\ell$ and suppose the image of
$G_{\mathbf Q(\mu_{\ell^\infty})}$ on $T_\ell$ has finite
index $m$ in $\operatorname{SL}_2(\mathbf Z_\ell)$.
At every projective level $\ell^r$, the scalar twist in
\textup{(D.32.9)} disappears. The acting image has index at
most $m$ in the ambient transitive group, so it has at most
$m$ orbits on each cusp fibre. There are at most $3m$
closed cusps and the rational unit rank is at most $3m-1$,
uniformly in $r$.
For the full cover, restricting further to $\ker\delta_c$
increases the image index by at most two. It gives the
uniform bounds $6m$ and $6m-1$ respectively.
The finite-index assumption is explicit; it has not been
deduced for arbitrary section systems.

\proofparagraph{Exact checks and the remaining scope.}
The standard-library script
\path{code/D_modular_cusp_orbits.py} enumerates the
$120$ matrices generated by \textup{(D.32.4)} modulo five,
checks the three inertia orders and left-coset right actions,
and verifies the three full and projective cusp counts.
It also enumerates $E_3(\mathbf F_7)$, checks the irreducible
Frobenius polynomial modulo five, generates all $480$
matrices of $\operatorname{GL}_2(\mathbf F_5)$ from that
Frobenius and a unipotent, and verifies
\textup{(D.32.14)} exactly over $\mathbf Q$.
Its output is \path{code/D_modular_cusp_orbits.json}.
These finite checks supplement the all-level proof.
The detailed source and route audit is
\path{audit/D_modular_level_unit_rank_obstruction.md}.
D.31's multiple-label problem and the finite-ramification
problem remain unresolved.
